\section{Step 1 Admit strides before touching samples}\label{h263-sad-s01}
Begin with the two callable shells in the supplied file. Keep the returned count
zero while checking arguments. Reject an unsigned stride below sixteen or above
4096. Apply the check independently to both inputs. A negative register pattern
is a large unsigned stride, not a request to walk rows backwards.

Use the same failure convention in scalar and MSA entries. A valid computation
that is still unfinished may return development status four; an argument rejected
by the completed interface returns one. Sharing the checks is reasonable, but it
is not required for an equivalent implementation.

\paragraph{Observe} In a prepared test copy, use stride fifteen. The function
must reject it without reading an input sample. Restore the worked stride 176
before continuing with the picture windows.
\Needspace{8\baselineskip}
\paragraph{Check} Why is stride fifteen invalid for this entry?
\begin{itemize}
\item[A.] The fixed comparison consumes sixteen active samples in every row.
\item[B.] Every MSA byte load requires an address divisible by sixteen.
\item[C.] A luma sample occupies two bytes.
\end{itemize}

\section{Step 2 Bound the complete address windows}\label{h263-sad-s02}
After the stride checks, compute $15s+16$ for each input in 64-bit registers.
The admitted stride bound makes these extent calculations small, but the input
pointer may still be near the end of the address space. Add the extent to its
start and use an unsigned comparison to detect a wrapping exclusive end.

Perform both checks before the first load from either source. Otherwise, a valid
first input could already have been read when the second pointer is rejected.
A mapped input window is a caller obligation; this arithmetic check does not
inspect the simulator's memory map or turn an arbitrary unmapped pointer into a
successful argument.

\paragraph{Observe} For stride 176, the minimum extent is 2656 bytes. For stride
sixteen it is 256. A top address whose addition wraps is rejected with count zero.
\Needspace{8\baselineskip}
\paragraph{Check} How much input does one successful comparison consume?
\begin{itemize}
\item[A.] Every byte in both contiguous extents.
\item[B.] Exactly 256 active sample bytes from each window.
\item[C.] Sixteen bytes from each whole picture.
\end{itemize}

\section{Step 3 Form one scalar difference}\label{h263-sad-s03}
Use separate row-start pointers and a count of sixteen rows. At the start of a
row, copy those starts into moving sample cursors and initialize a sixteen-sample
column count. Load one unsigned byte from each cursor. Their subtraction lies
between minus 255 and plus 255, so its absolute value fits an ordinary general
register without a special minimum-integer case.

Accumulate the magnitude in a wide sum. The current input's first worked sample
is 56 and the reference sample is 200, giving magnitude 144. Loading 200 with a
signed byte load would instead produce minus 56, so the wrong result can arise
before the absolute-value operation. Interpret the samples correctly at the load.

\paragraph{Observe} Stop immediately after forming the first magnitude. It is
144. No input byte has been modified.
\Needspace{8\baselineskip}
\paragraph{Check} Which loaded values should be subtracted for this pair?
\begin{itemize}
\item[A.] 56 and minus 56.
\item[B.] Minus 200 and 56.
\item[C.] 56 and 200.
\end{itemize}

\section{Step 4 Finish exactly sixteen scalar rows}\label{h263-sad-s04}
Advance both sample cursors by one after each pair, and repeat for sixteen pairs.
When a row finishes, advance the retained row starts by their respective strides.
The row starts and moving cursors have different jobs: adding the stride to a
cursor already sixteen bytes past the start would drift sixteen extra samples
on every row.

Finish sixteen rows and retain the complete sum. For the worked first row, eight
differences are 144 and eight are 132, giving 2208. Later rows overlap the bright
square differently, so multiplying 2208 by sixteen is not the block result.
The complete scalar sum is 7872. A row's neighbors and padding do not contribute.

\paragraph{Observe} The successful trace has 256 sample loads from each window.
Their addresses fit the sixteen active bytes in each row. The final sum is 7872.
\Needspace{8\baselineskip}
\paragraph{Check} Which pointer should receive the row stride?
\begin{itemize}
\item[A.] The pointer to the start of the current row.
\item[B.] The pointer already sixteen samples beyond that start, without adjustment.
\item[C.] The address of the returned sum.
\end{itemize}

\section{Step 5 Compare a whole row with MSA}\label{h263-sad-s05}
Start the MSA entry with the same admission checks and row count. Zero its wide
lane accumulator once for the whole block. For each row, load exactly sixteen
bytes from each active start and form unsigned byte difference magnitudes with
\texttt{asub\_u.b}. Keep the input samples and the byte differences conceptually
separate: their bit widths agree, but their meanings do not.

There is no scalar remainder within a sixteen-sample row. Do not load a second
vector to make an unaligned address convenient. A byte-form MSA load is permitted
at the supplied unaligned byte pointer; its entire sixteen-byte access must still
fit the active row.

\paragraph{Observe} The first worked row's magnitudes consist of eight 144s and
eight 132s by increasing sample address. All-zero against all-255 input produces
sixteen magnitudes of 255.
\Needspace{8\baselineskip}
\paragraph{Check} What is wrong with using \texttt{asub\_s.b} here?
\begin{itemize}
\item[A.] It always changes the memory byte order.
\item[B.] It interprets samples above 127 as negative values.
\item[C.] It returns a floating-point result.
\end{itemize}

\section{Step 6 Widen before accumulating}\label{h263-sad-s06}
Combine adjacent difference bytes into unsigned halfwords, then add those
halfwords to the block accumulator. For the first worked row, the adjacent pair
sums by increasing address are four 288s and four 264s. Each pair now needs more
than one byte even though its two original magnitudes individually fit.

The width proof applies to every admitted byte pattern, not just the selected
video. A halfword lane can receive at most $16\times510=8160$ over the entire
block. That bound permits wrapping halfword addition here without losing a carry.
It does not permit byte accumulation followed by a wider interpretation of its
already-wrapped result.

\paragraph{Observe} With all differences 255, each of the eight final halfword
lanes is 8160. Their total is 65\,280.
\Needspace{8\baselineskip}
\paragraph{Check} Which operation preserves the information needed for that total?
\begin{itemize}
\item[A.] Add rows in byte lanes, then reinterpret the bytes as halfwords.
\item[B.] Saturate every row total to 255.
\item[C.] Widen adjacent byte sums before accumulating across rows.
\end{itemize}

\section{Step 7 Advance rows without extending the reads}\label{h263-sad-s07}
After each vector row, advance the two row-start addresses by their own strides
and reduce the remaining-row count. Continue for sixteen rows. The loaded width
is sixteen bytes, but the worked row advance is 176 bytes. Those values agree
only for a tightly packed block copy, which this caller does not supply.

Do not clear the accumulator inside the loop. Equal inputs then hide the error,
because every row contributes zero. Use a maximum-difference block and the worked
moving-square block to expose lost earlier contributions. Keep the final row's
access bounded even if the next scratch pointer would never be dereferenced.

\paragraph{Observe} Each input generates sixteen full-width sample loads. No
load includes a seventeenth sample, a padding byte or the following row.
\Needspace{8\baselineskip}
\paragraph{Check} How many row-vector loads does this block need in total?
\begin{itemize}
\item[A.] Sixteen, shared between both inputs.
\item[B.] Thirty-two, sixteen from each input.
\item[C.] Two, one from each whole frame.
\end{itemize}

\section{Step 8 Reduce the accumulated lanes}\label{h263-sad-s08}
Reduce the eight halfwords into four words, then into two doublewords. Extract
both doublewords into general registers and add them. The final sum includes all
sixteen rows and all sixteen columns. Returning one doubleword alone would retain
only part of the block, even if that partial sum looks reasonable.

The two wide partials are bounded by 32\,640 each. A doubleword extraction cannot
turn either into a negative value. Keep the return count in a full-width general
register so its representation is unambiguous at the callable boundary. Compare
both target byte orders: a consistent permutation of byte positions can leave a
sum unchanged, while an incorrect extraction or missed half still fails.

\paragraph{Observe} Maximum-difference inputs return 65\,280. The worked windows
return 7872, and identical windows return zero.
\Needspace{8\baselineskip}
\paragraph{Check} What does the final scalar addition combine?
\begin{itemize}
\item[A.] Two partial sums covering the complete block together.
\item[B.] The addresses of the two source buffers.
\item[C.] Two competing motion vectors.
\end{itemize}

\section{Step 9 Return the result through the supplied caller}\label{h263-sad-s09}
Publish success status zero only after the whole block has been compared. Return
the SAD in \texttt{\$v1}, preserve saved registers and return through the caller's
address. A zero SAD is an ordinary successful result; it is not the failure
convention. Overlapping read-only windows are also legitimate arguments.

Prepare separate worked copies in scalar and MSA modes. Run each to completion.
The completed caller prints 7872 followed by a newline and exits successfully.
Before the functions are complete, its development message and failing process
status remain expected. Treat that message as progress feedback rather than a
completed codec result.

\paragraph{Observe} In Debug, stop before \texttt{sad\_caller\_return}'s first
instruction. Returned registers are zero and 7872 before the caller stores them
in \texttt{sad\_last\_status} and \texttt{sad\_last\_sum}.
\Needspace{8\baselineskip}
\paragraph{Check} What should identical valid windows return?
\begin{itemize}
\item[A.] Failure status one because no motion was found.
\item[B.] Development status four forever.
\item[C.] Success status zero and SAD zero.
\end{itemize}

\section{Step 10 Measure this boundary before adding search}\label{h263-sad-s10}
Count completed instructions from function entry through its return to a stopped
caller continuation. Include admission checks and reduction; exclude preparation
of the caller's data. Record which implementation and input produced the count.
An instruction count is reproducible simulator work, while host wall time also
includes scheduling and runtime effects.

The checked reference used 3245 completed instructions in scalar mode and 207 in
MSA mode for each admitted worked block. A second valid MSA implementation used
239, because it formed magnitudes with separate maximum, minimum and subtraction
operations. All produced the same exact sums and active-row accesses. Neither
207 nor the ratio between these counts is a grading threshold.

\paragraph{Observe} Record counts for equal, maximum-difference and worked
windows. Confirm correctness independently before reporting the counts.
\Needspace{8\baselineskip}
\paragraph{Check} What can these measurements establish by themselves?
\begin{itemize}
\item[A.] Completed instruction work for the named kernel boundaries and inputs.
\item[B.] The running time of the finished video encoder.
\item[C.] A requirement that every correct solution execute exactly 207 instructions.
\end{itemize}
